\documentclass[12pt]{article}
\usepackage{amsmath,amssymb,amsthm,geometry,algorithm,algpseudocode}
\geometry{a4paper,margin=1in}
\newtheorem{theorem}{Theorem}
\newtheorem{lemma}[theorem]{Lemma}
\newtheorem{corollary}[theorem]{Corollary}
\newtheorem{proposition}[theorem]{Proposition}
\theoremstyle{definition}
\newtheorem{definition}[theorem]{Definition}

\title{Project Euler Problem 456:\\Triangles Containing the Origin II}
\author{}
\date{}

\begin{document}
\maketitle

\section{Problem Statement}
Define:
\begin{align}
x_n &= (1248n \bmod 32323) - 16161 \\
y_n &= (8421n \bmod 30103) - 15051 \\
P_n &= \{(x_1, y_1), (x_2, y_2), \ldots, (x_n, y_n)\}
\end{align}

Let $C(n)$ be the number of triangles whose vertices are in $P_n$ which contain the origin in their interior. Find $C(2{,}000{,}000)$.

\section{Mathematical Analysis}

\subsection{Origin Containment via Cross Products}
A triangle with vertices $A$, $B$, $C$ contains the origin if and only if the cross products $A \times B$, $B \times C$, $C \times A$ all share the same sign, where $A \times B = a_x b_y - a_y b_x$.

\subsection{Angular Sweep Method}
For each point $P_i$, compute $\theta_i = \text{atan2}(y_i, x_i)$. Sort points by angle. A triangle does \textbf{not} contain the origin if and only if all three vertices lie within some open half-plane through the origin.

\begin{theorem}
Let $f(i)$ denote the number of points lying in the open semicircle $(\theta_i, \theta_i + \pi)$. Then:
\[
C(n) = \binom{n}{3} - \sum_{i=1}^{n} \binom{f(i)}{2}
\]
\end{theorem}

\begin{proof}
A triangle fails to contain the origin iff there exists a line through the origin with all three vertices on one side. Sweeping a half-plane, when the leading edge passes point $i$, point $i$ is the ``last'' point to enter. The points already inside form the set counted by $f(i)$. Each bad triangle is counted exactly once (by the point with the most clockwise position as boundary).
\end{proof}

\subsection{Two-Pointer Technique}
After sorting angles, maintain pointers $i$ and $j$ such that $\theta_j - \theta_i < \pi$. As $i$ advances, $j$ advances monotonically, giving $f(i) = j - i - 1$ in the circular arrangement.

\section{Complexity}
\begin{itemize}
\item Time: $O(n \log n)$ for sorting, $O(n)$ for sweep.
\item Space: $O(n)$.
\end{itemize}

\section{Answer}

\[
\boxed{333333208685971546}
\]

\end{document}
